\chapter{Weights and tangent measures}\label{sec:weighted-rectifiable-tangents}
An integrable weight can make a rectifiable surface measure locally finite even when its unweighted area is not. We prove the weighted tangent statement by exhausting the surface with subsets where the weight has a positive lower bound. This expands \cite{Sim18}, Chapter 3, Definition 1.7 and Theorem 1.9, printed pp.~90--92. Here $1\leq m<n$ and $U\subset\mathbb R^n$ is open. We write $\eta_{x,r}(y)=(y-x)/r$.

\begin{proposition}[Multiplying a unit-density tangent by an integrable weight]\label{prop:weighted-tangent-local-area}
Let $M\subset U$ be measurable, countably $m$-rectifiable and locally finite in $\mathcal H^m$ measure. Put $\sigma=\mathcal H^m\llcorner M$. For $f\in L^1_{\mathrm{loc}}(\sigma)$, at $\sigma$-almost every $x$ and for every $\phi\in C_c(\mathbb R^n)$,
\[
 \lim_{r\downarrow0}r^{-m}\int_M\phi((y-x)/r)f(y)\,d\mathcal H^m(y)
 =f(x)\int_{T_xM}\phi\,d\mathcal H^m.
\]
The weight may be signed. All integrals are local to $U$ for small enough $r$.
\end{proposition}
\begin{proof}
Choose a point where Theorem \ref{thm:rectifiable-forward-tangents}, Corollary \ref{cor:rectifiable-local-unit-density}, and the absolute-oscillation conclusion of Theorem \ref{thm:radon-function-differentiation} for $f$ all hold. Their simultaneous set has full $\sigma$ measure. For a test supported in $B_R$, the difference between the scaled integral with $f(y)$ and that with the constant $f(x)$ has absolute value at most
\[
 \|\phi\|_\infty r^{-m}\int_{B_{Rr}(x)}|f(y)-f(x)|\,d\sigma(y).
\]
The oscillation integral is $o(\sigma(B_{Rr}(x)))$, and the density limit gives $\sigma(B_{Rr}(x))=O(r^m)$. Thus the error tends to zero. The remaining constant-weight integral converges by the tangent theorem. The same full-measure set works for every compactly supported test, since the oscillation and density estimates hold at all small radii. A countable compact exhaustion proves the local-domain version.
\end{proof}

\begin{definition}[A tangent plane with respect to a positive weight]\label{def:weighted-approximate-tangent}
Let $M\subset U$ be $\mathcal H^m$-measurable, and let $\theta:M\to(0,\infty)$ be measurable and finite, with $\mu=\theta\mathcal H^m\llcorner M$ locally finite. An $m$-plane $P$ is a weighted approximate tangent at $x\in M$ if
\[
 \lim_{r\downarrow0}r^{-m}\int_M\phi((y-x)/r)\theta(y)\,d\mathcal H^m(y)
 =\theta(x)\int_P\phi\,d\mathcal H^m
 \quad(\phi\in C_c(\mathbb R^n)).
\]
This is the convention of \cite{Sim18}, Definition 1.7, p.~91. Positivity and finiteness of $\theta(x)$ make the plane unique when it exists. The definition does not require locally finite unweighted area.
\end{definition}

\begin{lemma}[Positive weights give finite-area pieces and the same null sets]\label{lem:positive-weight-area-exhaustion}
Under Definition \ref{def:weighted-approximate-tangent}, there are increasing measurable sets $M_j\subset M$ covering $M$, such that $M_j$ is bounded and has finite $\mathcal H^m$ measure, and $1/j\leq\theta\leq j$ on $M_j$. Further, $\mu(A)=0$ for a measurable $A\subset M$ if and only if $\mathcal H^m(A)=0$. Borel representatives can be chosen for $M$ and its restricted weighted measure.
\end{lemma}
\begin{proof}
Choose increasing compact sets $K_j\subset U$ whose union is $U$, and put $M_j=M\cap K_j\cap\{1/j\leq\theta\leq j\}$. They increase and cover $M$ because the weight is positive and finite. The bound $\mathcal H^m(M_j)\leq j\mu(K_j)<\infty$ proves finite area. For a measurable $A$, a zero Hausdorff measure gives a zero weighted integral; conversely $\mathcal H^m(A\cap M_j)\leq j\mu(A)$ and continuity from below show that a zero weighted integral forces $\mathcal H^m(A)=0$.

Each $M_j$ has a Borel hull with null Hausdorff difference, by the argument in Lemma \ref{lem:rectifiable-borel-representative}. Their union is a Borel representative of $M$ modulo an $\mathcal H^m$-null set. On this sigma-finite Hausdorff measure space, a measurable finite function has a Borel representative modulo a null set, by approximation with simple functions and Borel representatives of their level sets. Its weighted restriction is a locally finite Borel measure, hence Radon. Such changes do not change any integral or almost-everywhere assertion.
\end{proof}

\begin{lemma}[Recovering a tangent from a full-density restriction]\label{lem:tangent-full-density-restriction}
Let $\mu$ be Radon, $A$ Borel, $x\in A$, and suppose
\[
 \frac{(\mu\llcorner A)(B_r(x))}{\mu(B_r(x))}\to1,
 \qquad r^{-m}(\eta_{x,r})_\#(\mu\llcorner A)
 \rightharpoonup q\mathcal H^m\llcorner P
\]
for $q\in(0,\infty)$ and an $m$-plane $P$. Then $r^{-m}(\eta_{x,r})_\#\mu$ has the same limit, and
\[
 \mu(B_r(x)\setminus A)=o(r^m).
\]
\end{lemma}
\begin{proof}
The plane-tangent ball calculation gives $(\mu\llcorner A)(B_r)/(\omega_mr^m)\to q$. Dividing by the first assumed ratio gives $\mu(B_r)/(\omega_mr^m)\to q$. Subtracting the two expressions gives the stated small mass outside $A$. The difference of the two scaled integrals against a test supported in $B_R$ is bounded by its supremum norm times $r^{-m}\mu(B_{Rr}(x)\setminus A)$, which tends to zero. Thus the unrestricted tangent is the same. This conclusion needs the positive finite density of the restricted tangent as well as relative density one.
\end{proof}

\begin{theorem}[Weighted tangents without locally finite unweighted area]\label{thm:weighted-rectifiable-forward}
Suppose $M$ is countably $m$-rectifiable and satisfies Definition \ref{def:weighted-approximate-tangent}. Then the weighted approximate tangent exists at $\mathcal H^m$-almost every point of $M$, and
\[
 \lim_{r\downarrow0}\frac{\mu(B_r(x))}{\omega_mr^m}=\theta(x)
\]
at almost every such point. This proves the forward direction of \cite{Sim18}, Theorem 1.9, p.~91.
\end{theorem}
\begin{proof}
Use the Borel representatives and pieces $M_j$ in Lemma \ref{lem:positive-weight-area-exhaustion}. Each piece is rectifiable with finite unweighted area. Proposition \ref{prop:weighted-tangent-local-area}, applied to its restricted weight, gives the tangent $\theta(x)\mathcal H^m\llcorner T_xM_j$ for $\mu_j=\mu\llcorner M_j$ at $\mathcal H^m$-almost every $x\in M_j$. By Theorem \ref{thm:radon-function-differentiation} for $\chi_{M_j}$ with respect to $\mu$, one also has $\mu_j(B_r)/\mu(B_r)\to1$ at $\mu$-almost every point of $M_j$. The null-set equivalence makes both conclusions simultaneous outside an $\mathcal H^m$-null subset of $M_j$. Lemma \ref{lem:tangent-full-density-restriction} transfers the tangent to $\mu$. A countable union of exceptional subsets covers the only exclusions in $M$. The ball-density assertion follows either from that lemma or by continuous approximation of the unit-ball indicator. Positivity of $\theta$ is used when transferring relative density to absolute scale $r^m$.
\end{proof}

\begin{theorem}[Positive weighted plane tangents imply rectifiability]\label{thm:weighted-rectifiable-converse}
Under Definition \ref{def:weighted-approximate-tangent}, suppose the weighted approximate tangent exists at $\mathcal H^m$-almost every point of $M$. Then $M$ is countably $m$-rectifiable. Together with Theorem \ref{thm:weighted-rectifiable-forward}, this gives the weighted characterization in \cite{Sim18}, Theorem 1.9.
\end{theorem}
\begin{proof}
At a point with tangent $\theta(x)\mathcal H^m\llcorner P_x$, the continuous-cutoff and compact-cone arguments in Lemma \ref{lem:plane-tangent-density-cone} give positive lower density $\theta(x)\omega_m$ for $\mu$ and zero normalized mass in $C_{1/4}(x,P_x)\cap B_r(x)$. Apply Theorem \ref{thm:pointwise-plane-cone-cover} to the full-measure set of such points, using the Radon measure from Lemma \ref{lem:positive-weight-area-exhaustion}. The result is a countable graph cover; the omitted subset is Hausdorff-null by hypothesis. This proves rectifiability without asserting locally finite unweighted area. The positive density constant may vary with $x$, as allowed by the pointwise cone theorem.
\end{proof}

\begin{corollary}[The tangent plane is independent of a positive weight]\label{cor:weighted-tangent-independent}
Suppose a rectifiable $M$ carries two positive finite measurable weights $\theta,\rho$ whose weighted measures are locally finite. Their weighted approximate tangent planes agree at $\mathcal H^m$-almost every point of $M$, although their limiting plane measures have coefficients $\theta(x)$ and $\rho(x)$.
\end{corollary}
\begin{proof}
Exhaust $M$ by compactly located measurable pieces where both weights lie between $1/j$ and $j$. Each piece has finite unweighted area. On each piece Proposition \ref{prop:weighted-tangent-local-area} gives the same ordinary tangent plane for both weights, with their respective coefficients. Radon differentiation of the piece's indicator and Lemma \ref{lem:tangent-full-density-restriction} transfer each limit to its full weighted measure. Remove the countable union of exceptional Hausdorff-null subsets. The limiting planes coincide there because both were obtained from the same piece. This proof compares planes, not the unequal scalar multiplicities.
\end{proof}
